\title{xLSTM: Extended Long Short-Term Memory}

\begin{abstract}
During the 1990s, the Long Short-Term Memory (LSTM) framework was fundamentally shaped by the introduction of gating mechanisms and the constant error carousel, which quickly established LSTMs as a reliable building block for deep learning systems. These networks were instrumental in enabling the initial Large Language Models (LLMs). Nevertheless, the emergence of the Transformer paradigm, distinguished by its parallelizable self-attention mechanism, revolutionized sequence modeling and rapidly surpassed LSTM-based architectures, especially as model size increased. In light of these advances, we investigate the potential of LSTMs in contemporary language modeling by dramatically expanding their parameter count to the billion scale and incorporating innovative strategies derived from recent progress in LLMs, while explicitly addressing their traditional shortcomings. To this end, we propose the use of exponential gating alongside advanced normalization and stabilization procedures. Moreover, we reimagine the LSTM memory architecture, introducing: (i) sLSTM, which incorporates scalar-based memory and update mechanisms, paired with novel memory integration techniques, and (ii) mLSTM, characterized by fully parallelizable, matrix-form memory and a covariance-oriented update scheme. These new LSTM configurations are integrated into residual backbone frameworks, resulting in xLSTM blocks, which are then assembled in a stacked, residual fashion to form xLSTM architectures. Through the adoption of exponential gating and refined memory mechanisms, xLSTM models attain competitive, and sometimes superior, performance relative to the latest Transformer and State Space Model architectures, especially with respect to both accuracy and scalability.\\
Code available at: \url{https://github.com/NX-AI/xlstm}
\end{abstract}

\section{Introduction}
\label{sec:intro}

The Long Short-Term Memory (LSTM) ideas~\citep{Hochreiter:91,Hochreiter:97nips,Hochreiter:97}, i.e., the constant error carousel and gating, were introduced to overcome the vanishing gradient problem of recurrent neural networks~\citep{Hochreiter:91,Hochreiter:00book}:

\begin{align}
\label{eq:lstm_idea}
  c_t \ = \ f_t \ c_{t-1} \ + \ 
          i_t \  z_t \ , \quad  
          h_t \ = \ o_t \ \psi({c_t}) \ . 
\end{align}

The constant error carousel involves the cell state $c_{t-1}$ (depicted in green) being incremented by the cell inputs $z_t$, with the modification regulated via sigmoid gates (shown in blue). The update process is governed by the input gate $i_t$ alongside the forget gate $f_t$, whereas the output gate $o_t$ determines what the memory cell emits—specifically, the hidden state $h_t$. Before producing the hidden state, the cell state undergoes normalization or compression through $\psi$, after which the output gate modulates the resulting hidden state.

LSTMs have achieved significant success across a broad range of application areas \citep{Hochreiter:01,Hochreiter:07,Schmidhuber:15}, dominating the landscape of text generation until the emergence of Transformers in 2017~\citep{Vaswani:17}. Numerous sequence modeling challenges have showcased the capabilities of LSTMs, including, but not limited to, text production~\citep{Graves:13,Karpathy:15blog}, handwriting synthesis~\citep{Graves:13}, sequence-to-sequence tasks in machine translation~\citep{Sutskever:14nips}, program analysis~\citep{Zaremba:14arxiv}, generation of image descriptions~\citep{Karpathy:15,Hossain:19}, automated code generation~\citep{Karpathy:15blog}, hydrological simulations such as rainfall-runoff prediction~\citep{Kratzert:18,Kratzert:19}, and flood warning systems~\citep{Nearing:24}. In the context of reinforcement learning, LSTMs have demonstrated superior performance as sequence modeling tools in systems such as the AlphaStar agent for StarCraft II~\citep{Vinyals:17short}, OpenAI Five for Dota 2~\citep{OpenAI:19}, and controllers for magnetic confinement in nuclear fusion research~\citep{Degrave:22short}. The architectural strength of LSTMs lies in their capacity for abstraction, with a proven ability to capture complex semantic representations and maintain information across sequences within their cell states~\citep{Karpathy:15blog}. This phenomenon has been evidenced by the identification of number and syntax-specific neurons~\citep{Lakretz:19}, units specializing in linguistic attributes~\citep{Bau:19}, and neurons encoding sentiment~\citep{Radford:17arxiv}. Despite the advent of newer architectures, LSTMs remain integral to critical applications~\citep{Degrave:22short,Nearing:24} and continue to demonstrate enduring relevance.

Although LSTMs have achieved significant breakthroughs, they are subject to three principal drawbacks: (i) Irrevocable storage actions. This shortcoming is illustrated through the {\it Nearest Neighbor Search} task (refer also to Appendix~\ref{sec:appNearestNeighborSearch}), where, given a reference vector, a sequence must be traversed to locate the most similar vector and retrieve the corresponding value at the sequence’s conclusion. The mean squared error for this scenario appears in the left panel of Figure~\ref{fig:lstmProblems}, revealing that LSTM architectures find it difficult to update a previously stored value if a better match arises later in the sequence; by contrast, our xLSTM model leverages exponential gating to effectively address this issue. (ii) Constrained memory capacity, since all information has to be encoded within scalar cell states. We highlight this constraint using the {\it Rare Token Prediction} task. In the right panel of Figure~\ref{fig:lstmProblems}, one observes the token prediction perplexity from the Wikitext-103 dataset~\citep{Merity:17} segmented by token frequency. LSTMs encounter significant difficulties when predicting rare tokens, a direct consequence of limited storage expressivity. xLSTM confronts this with matrix-based memory to ameliorate performance. (iii) Inherent sequentiality imposed by hidden-to-hidden state transitions, which impedes parallel execution due to entanglement of memory across time steps.

These inherent constraints of LSTM models contributed to the rising popularity of Transformer architectures~\citep{Vaswani:17} in language modeling. This prompts the question: What are the achievable outcomes in language modeling once these barriers are addressed and LSTMs are scaled to match the capacities of present-day Large Language Models?

\section{Extended Long Short-Term Memory}
\label{sec:xLSTM}

In response to the challenges presented by the LSTM, Extended Long Short-Term Memory (xLSTM) introduces two significant changes to the foundational LSTM mechanism outlined in Equation~\eqref{eq:lstm_idea}. Specifically, xLSTM incorporates exponential gating and innovative memory arrangements, thereby broadening the LSTM framework with two variants: (i) sLSTM (detailed in Section~\ref{sec:sLSTM}), which operates with scalar memory, scalar updating, and memory mixing; and (ii) mLSTM (outlined in Section~\ref{sec:mLSTM}), characterized by a matrix-based memory and an update mechanism relying on the covariance (outer product), enabling complete parallelization. The introduction of exponential gating bolsters both sLSTM and mLSTM. To achieve parallel computation, mLSTM eliminates memory mixing, including hidden-to-hidden recurrent links. Additionally, both models can be generalized to multiple memory cells, allowing sLSTM to support inter-cell memory mixing. sLSTM is also adaptable to multiple heads, where memory mixing takes place solely between cells within each head, and not between different heads. This head-based design, combined with exponential gating, represents a novel scheme for mixing memory in sLSTM. In the case of mLSTM, handling multiple cells and multiple heads is functionally equivalent.

When these LSTM variants are incorporated into residual block modules, xLSTM blocks emerge (see Section~\ref{sec:xLSTMblock}). A network built by stacking these blocks in a residual manner forms an xLSTM architecture (refer to Section~\ref{sec:xLSTMarchitecure}). Figure~\ref{fig:xlstm_sketch} offers a visual overview of the xLSTM architecture and its constituent elements.

\subsection{Review of the Long Short-Term Memory}
\label{sec:orgLSTM}

The foundational concept of LSTM~\citep{Hochreiter:91,Hochreiter:97nips,Hochreiter:97} centers around the scalar memory cell, which serves both as a processing and a storage mechanism and mitigates the vanishing gradient issue~\citep{Hochreiter:91,Hochreiter:00book} through what is known as the constant error carousel—a mechanism realized via cell state updates. Three gates are integrated into this memory cell: an input gate, an output gate, and a forget gate, with the latter introduced by \citet{Gers:00}. 
At a given time step $t$, the LSTM memory cell is updated according to the following rules:

\begin{align}
c_t \ &= \  f_t \ c_{t-1} \ + \ i_t \ z_t &  & &\text{cell state} \\
h_t  \ &= \ o_t \ \tilde{h}_t \ , & \tilde{h}_t \ &= \ \psi \left( c_t\right)
  &\text{hidden state} \\
z_t \ &= \ \varphi \left( \tilde{z}_t \right) \ , 
  &\tilde{z}_t \ &=  \ \mathbf{w}^\top_{z} \ \mathbf{x}_t \ + \
  r_{z}  h_{t-1} \ + \  b_{z} \ \
  &\text{cell input} \\
i_t \ &= \ \sigma \left( \tilde{i}_t  \right) \ , 
  &\tilde{i}_t \ &= \ \mathbf{w}^\top_{i} \ \mathbf{x}_t \ + \
  r_{i}  \ h_{t-1} \ + \  b_{i} \ \
  &\text{input gate} \\
f_t \ &= \ \sigma \left(  \tilde{f}_t \right) \ , 
  &\tilde{f}_t \ &= \ \mathbf{w}^\top_{f} \ \mathbf{x}_t  \ + \
  r_{f}  \ h_{t-1} \ + \  b_{f} \ \
  &\text{forget gate} \\
o_t \ &= \ \sigma \left( \tilde{o}_t \right) \ , 
  &\tilde{o}_t  \ &= \ \mathbf{w}^\top_{o} \ \mathbf{x}_t \ + \
  r_{o}  \ h_{t-1} \ + \  b_{o} \ \
  &\text{output gate} 
\end{align}

The weight vectors $\mathbf{w}_{z}$, $\mathbf{w}_{i}$, $\mathbf{w}_{f}$, and $\mathbf{w}_{o}$ denote the parameters associated with connections from the inputs $\mathbf{x}_t$ to the cell input, the input gate, the forget gate, and the output gate, respectively. Similarly, the recurrent weights $r_{z}$, $r_{i}$, $r_{f}$, and $r_{o}$ represent the links from the previous hidden state $h_{t-1}$ to the cell input, input gate, forget gate, and output gate. The terms $b_{z}$, $b_{i}$, $b_{f}$, and $b_{o}$ stand for the respective biases. The activation functions $\varphi$ and $\psi$, often instantiated as $\tanh$, are applied to the cell input and hidden state, respectively. In this context, $\psi$ is instrumental in bounding the values of the cell state, preventing it from diverging. All gating mechanisms utilize the sigmoid activation, i.e., $\sigma \left( x \right)= 1/(1+\exp(-x))$. Later developments combined several scalar memory cells into a vector $\mathbf{c}_t \in \mathbb{R}^{d}$, enabling the application of recurrent weight matrices $\mathbf{R} \in \mathbb{R}^{d \times d}$, which facilitate interactions among memory cell outputs~\citep{Greff:15}; additional information can be found in Appendix~\ref{sec:apporgLSTM}. Studies analyzing model variations confirmed the necessity of all components within the memory cell architecture~\citep{Greff:15}.

\subsection{sLSTM}
\label{sec:sLSTM}

To enable LSTMs to modify their storage choices dynamically, we incorporate exponential gates (depicted in red), along with normalization and stabilization strategies. Specifically, we allow the input and forget gates to utilize exponential activation functions. For the purpose of normalization, we define a normalizer state that accumulates the products of the input gate and all subsequent forget gates.

The scalar sLSTM forward pass is:

\begin{align}
\label{eq:slstmforward}
c_t \ &= \  f_t \ c_{t-1} \ + \ i_t \ z_t & & 
 &\text{cell state} \\
n_t \ &= \  f_t \ n_{t-1} \ + \ i_t  & & 
 &\text{normalizer state} \\
h_t \ &= \ o_t \ \tilde{h}_t \ ,
  &\tilde{h}_t \ &= \ c_t / n_t
&\text{hidden state} \\
z_t \ &= \ \varphi \left( \tilde{z}_t \right) \ , 
  &\tilde{z}_t \ &=  \ \mathbf{w}^\top_{z} \ \mathbf{x}_t \ + \
  r_{z}  h_{t-1} \ + \  b_{z}
  &\text{cell input} \\
i_t \ &= \ \!\exp \left( \tilde{i}_t  \right) \ , 
  &\tilde{i}_t \ &= \ \mathbf{w}^\top_{i} \ \mathbf{x}_t \ + \
  r_{i}  \ h_{t-1} \ + \  b_{i}
  &\text{input gate} \\
f_t \ &= \ \sigma \left(  \tilde{f}_t \right) \ \text{OR} \
\!\exp \left(  \tilde{f}_t \right) \ ,
  &\tilde{f}_t \ &= \ \mathbf{w}^\top_{f} \ \mathbf{x}_t  \ + \
  r_{f}  \ h_{t-1} \ + \  b_{f}
  &\text{forget gate} \\
o_t \ &= \ \sigma \left( \tilde{o}_t \right) \ , 
  &\tilde{o}_t  \ &= \ \mathbf{w}^\top_{o} \ \mathbf{x}_t \ + \
  r_{o}  \ h_{t-1} \ + \  b_{o}
  &\text{output gate} 
\end{align}

We transfer the original LSTM gating techniques, i.e., input- and/or hidden-dependent gating plus bias term, to the new architectures. Exponential activation functions can lead to large values that cause overflows. Therefore, we stabilize gates with an additional state \mbox{$m_t$~\citep{Milakov:18arxiv}}:

\begin{align}
\label{eq:slstmstabil}
m_t \ &= \ \max \left( \log ( f_t ) + m_{t-1} , \log ( i_t ) \right) &\text{stabilizer state} \\
i'_t \ &= \ \!\exp \left( \log \left ( i_t \right) - m_t \right) \ = \ \!\exp \left( \tilde{i}_t - m_t \right) \ \, 
  &\text{stabil. input gate} \\
f'_t \ &= \ \!\exp \left( \log \left( f_t \right) + m_{t-1} - m_t \right) \ \, 
  &\text{stabil. forget gate}
\end{align}

In Appendix~\ref{sec:appsLSTM}, we demonstrate that substituting $f_t$ with $f'_t$ and $i_t$ with $i'_t$ during the forward pass leaves both the network’s output and the gradients of the loss function with respect to the parameters unaffected.

\paragraph{Novel Mechanisms for Memory Mixing.} Similar to the conventional LSTM, sLSTM allows for the presence of several memory cells (refer to Appendix~\ref{sec:appsLSTM}). Having multiple memory cells facilitates memory mixing by means of recurrent connections $\mathbf{R}_{\mathbf{z}}$, $\mathbf{R}_{\mathbf{i}}$, $\mathbf{R}_{\mathbf{f}}$, $\mathbf{R}_{\mathbf{o}}$, linking the hidden state vector $\mathbf{h}$ to memory cell input $\mathbf{z}$ and to the gate vectors $\mathbf{i}$, $\mathbf{f}$, and $\mathbf{o}$. Exponential gating brings a distinctive feature to memory mixing dynamics. The redesigned sLSTM architecture supports multiple heads, enabling memory mixing within individual heads, while disallowing interaction between separate heads. By incorporating heads and exponential gating in sLSTM, a novel strategy for mixing memory representations is established.

\subsection{mLSTM}
\label{sec:mLSTM}

To bolster the memory capabilities of LSTM models, we replace the traditional scalar cell $c \in \mathbb{R}$ with a matrix cell $\mathbf{C} \in \mathbb{R}^{d \times d}$. This allows memory access to be accomplished through matrix multiplication. At any time step $t$, the system is designed to encode a key-value vector pair—the key $\mathbf{k}_t \in \mathbb{R}^d$ and the value $\mathbf{v}_t \in \mathbb{R}^d$—following the nomenclature established in the Transformer architecture. Subsequently, at a later time $t+\tau$, it should be possible to recover the value $\mathbf{v}_t$ using a query vector $\mathbf{q}_{t+\tau} \in \mathbb{R}^d$. This scenario aligns with the framework of Bidirectional Associative Memories (BAMs) as described by \citep{Kohonen:72,Anderson:72,Nakano:72,Anderson:77}.

The covariance update rule~\citep{Sejnowski:77,Dayan:91} for storing a key-value pair is

\begin{align}
\mathbf{C}_t \ &= \ \mathbf{C}_{t-1} \ + \ \mathbf{v}_{t} \ \mathbf{k}_t^\top \ .
\end{align}

We apply layer normalization prior to the projection of inputs onto keys and values, ensuring these representations possess zero mean. The update mechanism for covariance aligns with the optimal criterion described in~\citep{Dayan:91}, which favors the highest possible separability among retrieved binary vectors, thus maximizing the signal-to-noise ratio. Enhanced separability can be achieved by restricting retrieval to pairwise associations and accepting the quadratic computational cost characteristic of attention mechanisms~\citep{Krotov:16,Krotov:17,Ramsauer:21}. Notably, this covariance update strategy corresponds with Fast Weight Programmers \citep{Schmidhuber:92ncfastweights,Schlag:21}, which have been further adapted by introducing a fixed decay factor applied to $\mathbf{C}_{t-1}$ and a constant learning rate for the term $\mathbf{v}_{t} \mathbf{k}_t ^\top$~\citep{Ba:16}. Drawing on this analogy, we embed the covariance update into the LSTM architecture: here, the forget gate is interpreted as the decay factor, the input gate as the learning rate, and the output gate as a modifier for the retrieved vector.

For this matrix memory, the normalizer state is the weighted sum of key vectors, where each key vector is weighted by the input gate and all future forget gates. Again, the normalizer state keeps record of the strength of the gates. Since the dot product between query and normalizer state can be close to zero, we use the absolute value of this dot product and lower bound it by a threshold (typically 1.0) as done previously~\citep{Sun:23arxiv}. The mLSTM forward pass is:

\begin{align}
\label{eq:mlstm_recurrent_begin}
\mathbf{C}_t \ &= \  f_t \ \mathbf{C}_{t-1} \ + \ 
  i_t \ \mathbf{v}_t \ \mathbf{k}_t^\top & &  &\text{cell state} \\
\mathbf{n}_t \ &= \  f_t \ \mathbf{n}_{t-1} \ + \ 
  i_t \ \mathbf{k}_t & &  &\text{normalizer state} \\
\mathbf{h}_t  \ &= \ \mathbf{o}_t \ \odot \ \tilde{\mathbf{h}}_t
\ , \qquad \qquad 
\tilde{\mathbf{h}}_t \ = \   \mathbf{C}_t \mathbf{q}_t \ / \ 
  \max \left\{ |\mathbf{n}_t^\top \mathbf{q}_t|, 1 \right\} 
   &&&\text{hidden state} \\
\mathbf{q}_t \ &= \ \mathbf{W}_q \ \mathbf{x}_t \ + \ \mathbf{b}_q  & &  &\text{query input} \\
\mathbf{k}_t \ &= \ \frac{1}{\sqrt{d}} \mathbf{W}_k \ \mathbf{x}_t \ + \ \mathbf{b}_k  & &  &\text{key input} \\
\mathbf{v}_t \ &= \ \mathbf{W}_v \ \mathbf{x}_t \ + \ \mathbf{b}_v  & &  &\text{value input} \\
i_t \ &= \ \!\exp \!  \! \left( \tilde{i}_t  \right) \ , 
 \qquad \qquad \ \ \ \ \,
  \tilde{i}_t \ = \ \mathbf{w}^\top_{i} \ \mathbf{x}_t \ + \  b_{i} \ \ \ 
  &&&\text{input gate} \\
f_t \ &= \ \sigma \!  \left(  \tilde{f}_t \right) \ \text{OR} \ \!\exp \!  \! \left(  \tilde{f}_t \right) , \ \ \: \!
\tilde{f}_t \ = \ \mathbf{w}^\top_{f} \ \mathbf{x}_t  \ + \
  b_{f}
  &&& \text{forget gate} \\
\label{eq:mlstm_recurrent_end}
\mathbf{o}_t \ &= \ \sigma \left( \tilde{\mathbf{o}}_t \right) \ , \qquad \qquad \quad \ \ \ \,
  \tilde{\mathbf{o}}_t  \ = \ \mathbf{W}_{\mathbf{o}} \ \mathbf{x}_t \ + \
  \mathbf{b}_{\mathbf{o}}
  &&&\text{output gate} 
\end{align}

The architecture of mLSTM supports several memory cells, similar to the conventional LSTM model. In mLSTM, the use of multiple heads functions identically to multiple cells due to the absence of memory mixing across them. To maintain stability in the exponential gating mechanisms of mLSTM, we apply stabilization methods identical to those employed in sLSTM (refer to Equation~\ref{eq:slstmstabil}). As memory components in mLSTM operate independently, the recurrence relation can alternatively be expressed in a parallelized form. Further elaboration is provided in Appendix~\ref{sec:appmLSTM}.

\subsection{xLSTM Architecture}
\label{sec:xLSTMarch}

\paragraph{xLSTM Blocks.}
\label{sec:xLSTMblock}
The role of an xLSTM block is to generate a nonlinear summary of preceding information within a high-dimensional embedding, thereby facilitating clearer discrimination among various sequences of past events or contextual settings. Accurate differentiation of these histories is essential for successful prediction of subsequent sequence elements, such as future tokens. Our design philosophy is guided by Cover’s Theorem~\citep{Cover:65}, which asserts that nonlinear embeddings in expanded dimensional spaces yield a greater likelihood for linear separability of patterns compared to their original spaces. We examine two types of residual block structures: (i) A residual block featuring a post up-projection (akin to architectures used in Transformers), which first produces a nonlinear summary in the initial space, follows with a linear mapping to an augmented dimensionality, applies a nonlinear activation, and then projects back linearly to the initial dimension; refer to the left panel of Figure~\ref{fig:Backbones} and the third column in Figure~\ref{fig:xlstm_sketch}. A comprehensive illustration of this structure can be found in Figure~\ref{fig:appsLSTM_detailed} in the appendix. (ii) A residual block incorporating a pre up-projection (in the manner of State Space Models), where the mapping to a higher-dimensional space is performed first,

The past information is aggregated in a non-linear fashion within the high-dimensional space, which is subsequently projected back into the original space using a linear transformation. In cases where an xLSTM block contains an sLSTM, the post up-projection block is employed. Conversely, if the xLSTM block houses an mLSTM, we apply the pre up-projection block to accommodate the increased memory capacity inherent to the higher-dimensional space. For further clarification, consult the left panel of Figure~\ref{fig:Backbones}, the third column of Figure~\ref{fig:xlstm_sketch}, or Figure~\ref{fig:appsLSTM_detailed} located in the appendix.

\paragraph{xLSTM Architecture.}
\label{sec:xLSTMarchitecure}
The xLSTM framework is developed by layering modular building blocks in a residual manner~\citep{Srivastava:15,He:16}. Our implementation follows the pre-LayerNorm~\citep{Ba:16b} residual architecture, a design prevalent in modern Large Language Models. Reference is provided in the final column of Figure~\ref{fig:xlstm_sketch}.

\subsection{Memory and Speed Considerations}

In contrast to Transformers, xLSTM architectures maintain linear computational costs and fixed memory usage with respect to the length of the sequence. The compressive nature of xLSTM's memory design makes it particularly advantageous for deployment in industrial scenarios and resource-limited edge devices.

Although the memory component in mLSTM does not introduce additional parameters, it relies on a computationally demanding $d \times d$ matrix for both storage and update operations. This design balances memory capacity against computational overhead. However, these operations can be executed in parallel on GPUs, significantly mitigating their impact on overall runtime.

Although mLSTM supports parallel processing much like FlashAttention~\citep{Dao:22,Dao:23} or GLA~\citep{Yang:23arxiv}, sLSTM is inherently sequential due to its memory mixing via hidden-to-hidden connections, which precludes parallelization. Nevertheless, we have engineered an efficient CUDA version with memory optimizations down to the register level for sLSTM, resulting in an execution speed that is typically less than twice as slow as mLSTM.

\section{Related Work}
\label{sec:related}

\paragraph{Linear Attention.} A number of strategies have been proposed to address the issue of quadratic computational complexity with respect to context length in the Transformer architecture, aiming to render attention operations linear in context length. The Synthesizer approach generates artificial attention weights independent of token-to-token correlations~\citep{Tay:20arxiv}. Linformer implements self-attention via a low-rank matrix and further achieves a linear approximation of this process~\citep{Wang:20arxiv}. Linear Transformer transforms the attention computation into a linear form~\citep{Katharopoulos:20}. The Performer uses a technique based on positive orthogonal random features to linearly approximate the attention softmax~\citep{Choromanski:21}. Meanwhile, alternatives such as Structured Global Convolution (SGConv)~\citep{Li:22} and Hyena Hierarchy~\citep{Poli:23} replace conventional attention mechanisms with accelerated long-range convolutions.

\paragraph{State Space Models.} State Space Models (SSMs) have garnered significant attention recently due to their linear scaling with sequence length and their competitive results relative to Transformers. Early work in this direction introduced the Structured State Space sequence model (S4)~\citep{Gu:21}, which was succeeded by a range of variants, such as the Diagonal State Space (DSS) model~\citep{Gupta:22}, Gated State Space (GSS) models~\citep{Mehta:22}, the S5 model~\citep{Smith:22}, Bidirectional Gated SSM (BiGS)~\citep{Wang:22}, H3 model~\citep{Fu:23}, and more recently, Mamba~\citep{Gu:24arxiv}.

\paragraph{Recurrent Neural Networks.} The use of Recurrent Neural Networks (RNNs) has been proposed as an alternative to Transformer and attention mechanisms, owing to their ability to scale linearly with respect to context length. Language modeling tasks have seen encouraging outcomes with RNNs utilizing Deep Linear Recurrent Units (LRUs)~\citep{Orvieto:23, De:24arxiv}, as well as with architectures such as the Hierarchically Gated Linear RNN (HGRN)~\citep{Qin:23} and its successor HGRN2~\citep{Qin:24arxiv}. RWKV~\citep{Peng:23arxivshort,Peng:24arxivshort}, a recognized RNN-based method for large-scale language modeling, has demonstrated performance levels comparable to Transformer models.

\paragraph{Gating.}
Gating, a foundational concept originating with LSTM, has been reimagined and utilized extensively across various recent neural architectures. Implementations of gating mechanisms appear in HGRN~\citep{Qin:23}, HGRN2~\citep{Qin:24arxiv}, Gated Linear Attention (GLA)~\citep{Yang:23arxiv}, Gated State Space (GSS) models~\citep{Mehta:22}, Bidirectional Gated SSM (BiGS)~\citep{Wang:22}, Moving Average Equipped Gated Attention (MEGA)~\citep{Ma:22}, as well as in RWKV~\citep{Peng:23arxivshort} and Mamba~\citep{Gu:24arxiv}.

\paragraph{Covariance Update Rule.} In order to boost the storage abilities of the architecture, we incorporated a matrix memory into the mLSTM cell, following a covariance-based update scheme. Strategies employing similar covariance update dynamics include Fast Weight Programmers~\citep{Schmidhuber:92ncfastweights,Schlag:21}, RWKV-5 and RWKV-6~\citep{Peng:24arxivshort}, Retention~\citep{Sun:23arxiv}, Linear Transformer~\citep{Katharopoulos:20}, and HGRN2~\citep{Qin:24arxiv}.

\paragraph{Most Related.} The models most similar in design to xLSTM are Retention~\citep{Sun:23arxiv}, RWKV~\citep{Peng:23arxivshort,Peng:24arxivshort}, and HGRN2~\citep{Qin:24arxiv}. All three incorporate either matrix memory and/or gating mechanisms. A key distinction, however, is that the newly proposed sLSTM permits memory mixing, unlike these prior methods. Memory mixing is crucial for addressing state tracking tasks, which enhances the expressive capacity of LSTMs beyond what is achievable with State Space Models (SSMs) and Transformers~\citep{Merrill:24,Deletang:23}. State tracking capabilities are essential, for instance, when interpreting code or monitoring entities throughout a lengthy text.

\paragraph{Residually Stacking Architectures.} The xLSTM architectures, akin to nearly all recent large-scale deep learning models, are designed by stacking modular components in a residual manner~\citep{Srivastava:15,He:16}.

This architectural approach laid the foundation for advances in deep convolutional networks~\citep{He:16} and enabled the development of Transformers~\citep{Vaswani:17}. Notably, Transformers are central to Large Language Models (LLMs) such as GPT-3~\citep{Brown:20short}, ChatGPT~\citep{Schulman:22short}, GPT-4~\citep{Achiam:23gpt4short}, Megatron-LM~\citep{Shoeybi:19}, Gopher~\citep{Rae:21short}, ERNIE 3.0 Titan~\citep{Wang:21arxivshort}, GLaM~\citep{Du:21glamshort}, Chinese M6~\citep{Lin:21}, AlexaTM 20B (multilingual)~\citep{Soltan:22}, OPT~\citep{Zhang:22opt}, Chinchilla~\citep{Hoffmann:22short}, BLOOM~\citep{Scao:22arxivshort}, GLM-130B~\citep{Zeng:22arxivshort}, LaMDA~\citep{Thoppilan:22short}, PaLM~\citep{Chowdhery:22short}, Llama~\citep{Touvron:23arxiv}, and Gemini~\citep{GeminiTeam:23arxiv,Reid:24arxiv}.

\section{Experiments}
\label{sec:Exp}

This section presents an empirical analysis of xLSTM, benchmarking it against existing approaches with an emphasis on language modeling tasks. Section~\ref{sec:ExpSyntethic} explores the distinct abilities of xLSTM through synthetic task evaluations. In Section~\ref{sec:ExpComparison}, we examine validation perplexity across contemporary language modeling techniques, all trained using 15B tokens from SlimPajama~\citep{Soboleva:23}, and carry out ablation experiments for xLSTM on the same data. Additionally, we investigate the scaling patterns of each method in line with prior work by \citet{Kaplan:20} and \citet{Brown:20short}.

Section~\ref{sec:ExpLanguage} extends the evaluation to a comprehensive language modeling study. Here, xLSTM is juxtaposed with the top methods from Section~\ref{sec:ExpComparison}, with training performed on 300B tokens sourced from SlimPajama~\citep{Soboleva:23}. The investigation proceeds in four steps: we first analyze performance in handling extended context lengths, followed by comparisons of validation perplexity and results on downstream tasks~\citep{Sutawika:24}; next, we review outcomes across 571 textual domains from the PALOMA language benchmark~\citep{Magnusson:23arxivshort}; and finally, we evaluate scaling properties of all techniques given a twenty-fold increase in training data.

\label{sec:xlstm_blocks_notation}
Throughout our experiments, we denote the ratio of mLSTM- to sLSTM-based xLSTM blocks using the format xLSTM[$a$:$b$], where $a/b$ corresponds to the proportion of each type. For instance, xLSTM[7:1] signifies that in a sequence of eight blocks, seven utilize mLSTM and one uses sLSTM. In scenarios involving a total of 48 blocks—a typical setting—this specification results in 42 mLSTM-based and 6 sLSTM-based blocks. Additionally, all model variants include pre up-projection layers for mLSTM blocks and post up-projection layers for sLSTM blocks in every experiment.

\subsection{Synthetic Tasks and Long Range Arena}
\label{sec:ExpSyntethic}
To begin, we evaluate how well the exponential gating mechanism with memory mixing in xLSTM performs on formal language tasks~\citep{Deletang:23}. Next, we investigate the impact of xLSTM's matrix memory on the Multi-Query Associative Recall challenge~\citep{Arora:23arxiv}. Lastly, we measure xLSTM's ability to handle lengthy sequences using the Long Range Arena benchmark~\citep{Tay:21}.

\paragraph{Evaluation of xLSTM's Exponential Gating and Memory Mixing Capabilities.} We evaluate the effectiveness of xLSTM's novel exponential gating mechanism combined with memory mixing, a feature theorized to facilitate the resolution of state tracking challenges~\citep{Merrill:24, Merrill:23}. To support assessment of length generalization, we adapt and enhance the formal language benchmark suite introduced by \citet{Deletang:23} to permit training on sequences of varying lengths. More comprehensive task specifications and additional results are provided in Appendix~\ref{sec:appExpSynthetic-formalLang}. 

Our analysis juxtaposes xLSTM against a range of architectures, such as Transformers, State Space Models, and different classes of Recurrent Neural Networks. For each method, we report task-specific accuracy, restricting measurement to the pertinent tokens. Accuracies are normalized from 0 (chance) to 1 (ideal). We conduct comparisons between the 2-block versions of the respective models: xLSTM[0:1] (consisting solely of sLSTM), xLSTM[1:0] (only mLSTM), xLSTM[1:1], Llama, Mamba, RWKV, Retention, Hyena, LSTM, and LSTM incorporated within Transformer blocks (LSTM (Block)). Experimental outcomes are depicted in Figure~\ref{fig:formal-main}.

Architectures that lack memory mixing and, by extension, state tracking—such as Transformers and State Space Models—are unable to address problems like those defined by regular grammars (e.g., the parity task). This observation aligns with prior research, which indicates that both Transformers and State Space Models possess inherent limitations relative to the expressive capacity of RNNs~\citep{Merrill:24, Merrill:23, Deletang:23}.

\paragraph{Assessment of xLSTM Memory Capacity Using Associative Recall Tasks.} In this study, we evaluate the newly introduced matrix memory in xLSTM by measuring its performance on the Multi-Query Associative Recall benchmark~\citep{Arora:23arxiv}. This task requires each sequence to encode and subsequently recall a set of randomly sampled key-value pairs selected from an extensive vocabulary. To increase the challenge beyond the original formulation, we raise the upper limit of key-value pairs to 256 and expand the context window to 2048 tokens, thereby providing a more stringent test of model memory capacities. We benchmark 2-block configurations of Llama, Mamba, RWKV-5, RWKV-6, xLSTM[1:1], and xLSTM[1:0], with evaluation based on their accuracy in recalling the correct pairs. Given that memory in Transformer models like Llama scales exponentially with the coding dimension~\citep{Ramsauer:21}, they are considered the performance benchmark for this evaluation. Results are reported in Figure~\ref{fig:mqar-main}. 
Among non-Transformer contenders, xLSTM[1:1] consistently achieves superior results, including for smaller architectures. Notably, the inclusion of the sLSTM block does not compromise but rather augments the model's memory, a trend that is particularly noticeable on the hardest task variant with 256 pairs. Supplementary findings are included in Appendix~\ref{sec:appExpSynthetic-mqar}, where extrapolation experiments suggest that xLSTM's augmented memory is capable of generalizing to contexts extending beyond those encountered during training.

\paragraph{Test of xLSTM's Long Context Capabilities on Long Range Arena.} In order to evaluate how xLSTM manages lengthy sequences and extensive contexts, we conduct a comparative analysis of various approaches on the Long Range Arena~\citep{Tay:21}. Across all tasks, xLSTM achieves robust results, indicating that its architectural design is highly effective for diverse long-context challenges.

For more details, see Appendix~\ref{sec:appExpSynthetic-lra}.

\subsection{Method Comparison and Ablation Study}
\label{sec:ExpComparison}

In order to investigate the central question of our study—specifically, the capabilities of our novel LSTM variants when applied to large-scale language modeling—we conduct experiments where xLSTMs, Transformers, State Space Models, and additional baselines are all trained on 15B tokens extracted from the SlimPajama dataset, using an identical auto-regressive training paradigm. After training, these models are evaluated on the validation set, and ablation analyses are carried out on the xLSTMs.

\paragraph{Comparison of xLSTM with Existing Approaches.} To facilitate a fair evaluation, we trained all models on a dataset comprising 15B tokens from SlimPajama~\citep{Soboleva:23}. Model performance was assessed using perplexity scores on a held-out validation set. The following models were compared: xLSTM (introduced in this work), GPT-3 (Transformer architecture)~\citep{Brown:20short}, Llama (Transformer architecture)~\citep{Touvron:23arxiv}, H3 (SSM)~\citep{Fu:23}, Mamba (SSM)~\citep{Gu:24arxiv}, RWKV-4 (RNN)~\citep{Peng:23arxivshort}, RWKV-5 (RNN)~\citep{Peng:24arxivshort}, RWKV-6 (RNN)~\citep{Peng:24arxivshort}, GLA (linear Transformer)~\citep{Yang:23arxiv}, HGRN (RNN)~\citep{Qin:23}, HGRN2 (RNN)~\citep{Qin:24arxiv}, RetNet (linear Transformer)~\citep{Sun:23arxiv}, Hyena (linear Transformer)~\citep{Poli:23}, and variants xLSTM[1:0], xLSTM[7:1] as described in Section~\ref{sec:xlstm_blocks_notation}. 
All models employed mixed precision training; however, for RWKV-5, RWKV-6, GLA, and HGRN2, the process did not use PyTorch's automated mixed precision (further information is available in Appendix Section~\ref{sec:appGenTrainProc}).
We grouped methods as follows: (a) Transformers, (b) State Space Models (SSMs), and (c) Recurrent Neural Networks (RNNs) plus linear Transformers. Linear Transformers are characterized by the replacement of the original Transformer attention with linear alternatives. Each model was configured to match a GPT-3 baseline containing 350M parameters, with an embedding dimension of 1024 and 24 residual blocks. Notably, only GPT-3 employs weight sharing between token embeddings and output embeddings, resulting in a parameter count reduction. Results presented in Table~\ref{tab:spaj15b_model_results} demonstrate that xLSTM achieves superior perplexity on the validation set compared to all other methods evaluated. For supplementary details, refer to Appendix~\ref{sec:appExpComparison}.
Scaling behavior illustrated in Figure~\ref{fig:temp_sclaw15B} suggests that xLSTM maintains its advantage as model size increases.

\paragraph{Ablation Studies.}
As presented in Table~\ref{tab:spaj15b_model_results} and Figure~\ref{fig:temp_sclaw15B}, xLSTM demonstrates superior performance in language modeling tasks after being trained on 15B tokens sourced from SlimPajama. To systematically analyze the architectural modifications from LSTM to xLSTM, we incrementally transform a standard LSTM architecture into its xLSTM counterpart. The initial step involves embedding LSTM layers within a pre-LayerNorm residual backbone. In the next phase, a post up-projection module is incorporated. The final stage completes the transition with the addition of exponential gating and matrix memory mechanisms.

The results are shown in Table~\ref{tab:ablstudies} (top). 

Through ablation studies, the substantial enhancement in performance is linked to the incorporation of both exponential gating and matrix memory. Given the critical role of gating in both RNNs and State Space Models, various gating strategies were examined in further ablations. The results, presented in Table~\ref{tab:ablstudies} (bottom), indicate that making each gate trainable and dependent on the input leads to cumulative improvements. Further analysis regarding specific backbone elements can be found in Appendix~\ref{sec:appAblDetails}.

\subsection{xLSTM as Large Language Model}
\label{sec:ExpLanguage}

In this work, we extend our investigation to comprehensive language modeling tasks by evaluating the capabilities of xLSTM as a large language model (LLM). To facilitate this, we scale up the training dataset, using 300B tokens from SlimPajama—matching the token counts employed in models such as Mamba~\citep{Gu:24arxiv} and Griffin~\citep{De:24arxiv}.

Our analysis involves benchmarking xLSTM against RWKV-4, Llama, and Mamba, each representing their respective methodological categories delineated in Section~\ref{sec:ExpComparison}. RWKV-4 serves as the RNN baseline, since suitable precision settings for RWKV-5, RWKV-6, and HGRN2 were established only after the commencement of the 300B token training sessions (refer to Appendix~\ref{sec:appGenTrainProc} for details).

We train models at multiple scales—125M, 350M, 760M, and 1.3B parameters—and systematically probe their ability to extrapolate to longer sequences. Model performances are evaluated on validation data, and we further measure their effectiveness on downstream applications, their language modeling proficiency across 571 textual domains within the PALOMA benchmark, and their compliance with established scaling laws.

\paragraph{Sequence Length Extrapolation.}
\label{sec:spaj300B_ctx_extrapolation}
To begin, we investigate the ability of large-scale (1.3B parameter) variants of xLSTM, RWKV-4, Llama, and Mamba to extrapolate to longer sequence lengths. Each model is initially trained with a context window of 2048 tokens and subsequently evaluated with sequence lengths extending up to 16384 tokens. The outcomes are depicted in Figure~\ref{fig:spaj300B_seq_extrapolation}. Unlike alternative architectures, xLSTM sustains low perplexity even as the input context grows.

\paragraph{Validation Perplexity and Downstream Tasks.}
\label{sec:spaj_downstream_eval}
Next, across various model scales, we assess xLSTM, RWKV-4, Llama, and Mamba by measuring validation perplexity on the SlimPajama dataset for the next token prediction task as well as on additional benchmarks that probe for commonsense reasoning capabilities.

The third column in Table~\ref{tab:lmeval} displays the perplexity scores on the validation set achieved by various approaches. Among all model sizes, both xLSTM[1:0] and xLSTM[7:1] yield the lowest perplexity, indicating their superiority. The remaining columns in Table~\ref{tab:lmeval} summarize outcomes on downstream tasks. Across nearly all tasks and model configurations, xLSTM outperforms other methods; Mamba surpasses xLSTM only in selected instances within the ARC task. More comprehensive information can be found in Appendix~\ref{sec:appExpLanguage}.

\paragraph{Performance on PALOMA Language Tasks.} Furthermore, across all parameter sizes, we evaluate the effectiveness of xLSTM, RWKV-4, Llama, and Mamba models in predicting the subsequent token on PALOMA language tasks~\citep{Magnusson:23arxivshort}. The metric used for assessment is perplexity, calculated for next token prediction across 571 diverse text domains, covering sources such as nytimes.com and r/depression on Reddit.

As summarized in Table~\ref{tab:ppleval}, perplexity scores are presented in two sections: general language modeling benchmarks (the first seven columns) and more specific domain tasks (the last five columns). The results demonstrate that xLSTM[1:0] yields superior performance to xLSTM[7:1] in these settings.

In comparison, xLSTM[1:0] achieves lower perplexity than Mamba on 568 of 571 text domains (99.5\%), surpasses Llama on 486 out of 571 domains (85.1\%), and outperforms RWKV-4 in 570 out of 571 domains (99.8\%). Further information is provided in Appendix~\ref{sec:appExpLanguage}.

\paragraph{Scaling Laws.} Next, we examine the power-law scaling trend, which facilitates the prediction of performance outcomes for increased model sizes~\citep{Kaplan:20,Brown:20short}. Figure~\ref{fig:temp_sclaw300B} illustrates these scaling trends. Although all models demonstrate analogous scaling patterns, they differ in their intercepts. Among the group, RWKV-4 achieves the lowest performance, with Llama and Mamba following behind. The xLSTM model consistently outperforms Mamba, showing a comparative advantage similar to that between Mamba and Llama. This scaling pattern suggests that as model size grows, xLSTM is likely to maintain its competitive performance against both Transformer and State-Space architectures.

\textbf{Generation Times and Maximal Throughput.} In Figure~\ref{fig:inference_time_generation}, we present the results of our evaluation of text generation latency, and we report the corresponding maximum throughput in Figure~\ref{fig:inference_time_decoding_throughput} (left) for the xLSTM models at the 1.3B parameter scale. We benchmark these models against comparably-sized implementations of Mamba, Llama, and RWKV available from HuggingFace, with the Llama models leveraging a static key-value cache. During our evaluation, we observed that it was not possible to use full cache compilation for the Transformer or compile the Mamba model with \texttt{torch.compile}. All generation benchmarks were conducted with a batch size of 1 and a pre-fill length of 16 tokens, a scenario intentionally chosen to favor Transformer architectures.

The data in Figure~\ref{fig:inference_time_generation} reveals that xLSTM and the other recurrent architectures (namely, Mamba and RWKV-4) exhibit linear complexity in generation time, contrasting with the quadratic complexity manifested by the Llama model. For the throughput comparisons, we experimented with a range of batch sizes and pre-fill values for the Llama baselines. As depicted in Figure~\ref{fig:inference_time_decoding_throughput} (right), the constant memory usage of xLSTM supports substantially larger batch sizes than Llama, resulting in xLSTM achieving superior throughput.

\section{Limitations}

(i) Unlike mLSTM, the memory mixing mechanism in sLSTM inherently restricts parallelizability, impeding rapid parallel execution. However, we engineered a custom CUDA kernel for sLSTM, which currently operates at less than double the execution time compared to our parallelized mLSTM implementation. (ii) Existing CUDA kernels supporting mLSTM lack optimization, resulting in performance that is approximately four times slower relative to FlashAttention or the scan technique used in Mamba. Enhanced CUDA kernel speeds are likely achievable by leveraging approaches similar to FlashAttention. (iii) The computational demands of mLSTM’s matrix memory are elevated because it processes $d \times d$ matrices. Nevertheless, since parameter-free memory operations can be parallelized via conventional matrix procedures, the incremental runtime attributed to this complexity is relatively negligible. (iv) Proper initialization of the forget gates is essential. (v) Because the matrix memory does not depend on the sequence length, increasing the latter may result in memory overload when dealing with extensive contexts. Yet, this has not posed a significant issue for contexts as long as 16k; refer to Section~\ref{sec:spaj300B_ctx_extrapolation} for details. (vi) Owing to the considerable computational intensity of large language modeling experiments, neither the architecture nor the hyperparameters—particularly for larger xLSTM models—have been fully optimized. We expect that substantial architectural and hyperparameter tuning will be required to maximize xLSTM’s capabilities.

\section{Conclusion}
Our investigation has addressed the central question: To what extent can LSTM-based language modeling be advanced by scaling up to billions of parameters? At present, the evidence suggests: ``As far as contemporary approaches such as Transformers or State Space Models''. In this work, we introduced xLSTM, augmenting the original LSTM with exponential gating, memory mixing, and an innovative memory configuration. xLSTM demonstrates strong performance in language modeling benchmarks, rivaling established methods like Transformers and State Space Models. Observed scaling laws imply that more expansive xLSTM models may become prominent alternatives to current Large Language Models founded on Transformer architectures. Furthermore, xLSTM holds promise for significant contributions in domains including Reinforcement Learning, Time Series Prediction, and physical system modeling.

\end{document}